\documentclass[11pt, a4paper]{article}

\usepackage[T1]{fontenc}
\usepackage[utf8]{inputenc}
\usepackage{tgpagella}
\usepackage[spanish, es-noshorthands]{babel}
\usepackage{amsmath, amssymb, bm}
\usepackage{tabularx}
\usepackage{booktabs}
\usepackage[most]{tcolorbox}
\usepackage[margin=2.5cm]{geometry}
\usepackage{ragged2e}
\usepackage{fancyhdr}

\pagestyle{fancy}
\fancyhf{}
\setlength{\headheight}{16pt}
\lhead{\textbf{UCB Sede Tarija} | Ing. Mecatrónica}
\chead{}
\rhead{IMT-121 Circuitos Electrónicos I | Observation Checklist}
\lfoot{Prof: M.Sc. Ing. Bernardo Quiroga Turdera}
\rfoot{Página \thepage}
\renewcommand{\headrulewidth}{0.4pt}

\definecolor{ucbblue}{RGB}{0, 51, 102}

\begin{document}

\begin{tcolorbox}[colback=ucbblue!5!white, colframe=ucbblue, title=\textbf{Lista de Observación: Validaciones en Simulación Transitoria}]
    \centering \textbf{Registro de Desempeño Operativo en el Laboratorio Computacional}
\end{tcolorbox}

\noindent \textit{Instrucciones para el Instructor: Marque objetivamente la capacidad del estudiante / equipo para transicionar de la predicción analítica a la ejecución e interpretación computacional de LTspice.}

\vspace{0.4cm}
\textbf{Estudiante / Grupo Evaluado:} \hrulefill \quad \textbf{Fecha:} \hrulefill

\vspace{0.4cm}
\renewcommand{\arraystretch}{1.6}
\noindent
\begin{tabularx}{\textwidth}{>{\RaggedRight\arraybackslash}X >{\centering\arraybackslash}p{1.2cm} >{\centering\arraybackslash}p{1.8cm} >{\centering\arraybackslash}p{1.2cm} >{\RaggedRight\arraybackslash}p{3cm}}
\toprule
\textbf{Evidencia Observable} & \textbf{Sí} & \textbf{Con Apoyo} & \textbf{No} & \textbf{Nota Descriptiva} \\ 
\midrule
Predijo analíticamente $v_C(t_s^+)$ ANTES de ejecutar el simulador & $\square$ & $\square$ & $\square$ & \\ \addlinespace
Calculó la asíntota DC $v_C(\infty)$ aislando el puerto & $\square$ & $\square$ & $\square$ & \\ \addlinespace
Halló $R_{eq}$ apagando debidamente la fuente $V_s$ & $\square$ & $\square$ & $\square$ & \\ \addlinespace
Calculó e interpretó físicamente la escala $\tau$ & $\square$ & $\square$ & $\square$ & \\ \addlinespace
Identificó correctamente el flanco en la fuente \texttt{PULSE} & $\square$ & $\square$ & $\square$ & \\ \addlinespace
Configuró una directiva de simulación \texttt{.tran} coherente & $\square$ & $\square$ & $\square$ & \\ \addlinespace
Ejecutó o reprodujo el software de manera verificable & $\square$ & $\square$ & $\square$ & \\ \addlinespace
Graficó \texttt{V(a)} distinguiendo superposición de \texttt{V(in)} & $\square$ & $\square$ & $\square$ & \\ \addlinespace
Pudo leer el curso en $t_s + \tau$ en vez de $t=\tau$ & $\square$ & $\square$ & $\square$ & \\ \addlinespace
Justificó analíticamente si observó una ligera discrepancia & $\square$ & $\square$ & $\square$ & \\ \addlinespace
Predijo correctamente el efecto temporal de duplicar $C$ & $\square$ & $\square$ & $\square$ & \\ 
\bottomrule
\end{tabularx}

\vspace{0.6cm}
\begin{tcolorbox}[colback=gray!5, colframe=gray!50, title=\textbf{Clasificación Formativa de Estado Operativo (Sin Valor Sumativo)}]
Marque la etiqueta que mejor describa el estado general del estudiante tras la sesión:
\begin{itemize}\setlength\itemsep{0.2em}
    \item[$\square$] \texttt{ANALYTIC\_PREDICTION\_READY}: El modelo analítico es sólido, pero la parte de uso del software falló o quedó incompleta.
    \item[$\square$] \texttt{TRANS\_CONFIGURATION\_SUPPORTED}: Necesita asistencia para comprender qué significan los parámetros del bloque \texttt{PULSE} o del \texttt{.tran}.
    \item[$\square$] \texttt{INDEPENDENT\_EXECUTION\_VERIFIED}: Capaz de cerrar el bucle Teoría-SPICE por sí mismo.
    \item[$\square$] \texttt{INTERPRETATION\_READY}: Las gráficas están, pero requiere refuerzo para aprender a extraer datos y conclusiones usando cursores.
    \item[$\square$] \texttt{ACCESS\_FOLLOW\_UP}: El equipo o software limitó/impidió la ejecución y debe programarse un seguimiento posterior.
\end{itemize}
\end{tcolorbox}

\end{document}
